\documentclass[11pt,letterpaper]{article}

% ---------- Typography: matches the Week 1 proposal document ----------
\usepackage[margin=1.05in,top=0.95in,bottom=1in]{geometry}
\usepackage[T1]{fontenc}
\usepackage{mathpazo}            % Palatino body
\usepackage{helvet}              % Helvetica for headings and labels
\usepackage{microtype}
\usepackage{xcolor}
\usepackage[colorlinks=true,urlcolor=navy,linkcolor=navy]{hyperref}
\urlstyle{sf}
\usepackage{enumitem}
\usepackage{fancyhdr}
\usepackage{lastpage}

\definecolor{navy}{HTML}{1F3864}
\definecolor{midgray}{HTML}{595959}
\definecolor{hairline}{HTML}{C8C8C8}

\linespread{1.10}
\setlength{\parindent}{0pt}
\setlength{\parskip}{3.5pt}

\pagestyle{fancy}
\fancyhf{}
\renewcommand{\headrulewidth}{0pt}
\fancyfoot[L]{\sffamily\fontsize{7.5}{9}\selectfont\color{midgray}ENNACIRI \textbar\ LLM BRAINSTORM ACTIVITY \textbar\ JULY 2026}
\fancyfoot[R]{\sffamily\fontsize{7.5}{9}\selectfont\color{midgray}\thepage\ OF \pageref{LastPage}}

\newcommand{\sectionhead}[1]{%
  \vspace{6pt}%
  {\sffamily\bfseries\fontsize{9.5}{12}\selectfont\color{navy}\MakeUppercase{#1}}\par
  \vspace{-9pt}\textcolor{hairline}{\rule{\textwidth}{0.5pt}}\par\vspace{0pt}}

\newcommand{\answer}[1]{#1\par\vspace{3pt}}

\begin{document}

% ---------- Title block ----------
\vspace*{-14pt}
{\sffamily\fontsize{8.5}{11}\selectfont\color{midgray}MSDS 696 DATA SCIENCE PRACTICUM II \quad\textbar\quad WEEK 1 LLM BRAINSTORM ACTIVITY}\par\vspace{10pt}
{\sffamily\bfseries\fontsize{21}{25}\selectfont\color{navy}LLM Brainstorm Activity}\par\vspace{6pt}
{\fontsize{12}{16}\selectfont\itshape\color{midgray}How I used an LLM to go from topic ideas to a defensible problem}\par\vspace{9pt}
{\sffamily\fontsize{9}{12}\selectfont Oussama Ennaciri \quad\textbar\quad July 6, 2026}\par\vspace{5pt}
\textcolor{navy}{\rule{\textwidth}{1.2pt}}

% ---------- The exchange ----------
\sectionhead{The exchange}
\answer{I brainstormed with Claude (Claude Code) over several sessions: first to understand the course requirements and the curated project list, then to explore topic ideas from my own interests, and finally to pressure-test the problem statement piece by piece. The full exchange is long; the excerpts submitted with this activity show the turning points, including the one where I caught and corrected the LLM's framing of the prediction window.}

% ---------- The three lines ----------
\sectionhead{What it suggested}
\answer{Project directions built from my interests (a flight-delay insurance model, bank distress prediction), the Prompt Corrective Action ``undercapitalized'' tier as an objective target label, the FDIC BankFind API as the data source, and a ``within four quarters'' prediction window.}

\sectionhead{What I kept}
\answer{The bank-distress early-warning idea, because it fits my broker-dealer background; the PCA-based target; the FDIC API as the primary source; and the Silicon Valley Bank framing for the executive summary.}

\sectionhead{What I rejected and why}
\answer{The ``within four quarters'' window as first proposed: I caught that a four-quarter cumulative window can mean as little as one quarter of actual warning, so we reframed the question to ``at risk, and how soon?'' I also rejected Kaggle mirrors in favor of primary government sources, and rewrote the pitch and personal angle because the LLM's versions were too dramatic.}

\end{document}
